\section*{Capítulo \ref{ch:g9:alternating-current} --- La corriente alterna}

\begin{solution}{exo:g9:alternating-current:1}
La CC circula siempre en un sentido bajo una \omterm{def:g8:voltage-voltmeter:voltage}{tensión} constante ---
el régimen de la \omterm{def:g3:first-electric-circuit:battery}{pila}. La CA invierte el sentido periódicamente
bajo una \omterm{def:g8:voltage-voltmeter:voltage}{tensión} que oscila --- el régimen del enchufe.
\end{solution}

\begin{solution}{exo:g9:alternating-current:2}
$T$: la \omterm{def:g2:measuring-time:duration}{duración} de un ciclo completo; $f$: los ciclos por
segundo; $f = 1/T$ y $T = 1/f$. Para \qty{50}{Hz}: $T =
\qty{0.02}{s}$. Para $T = \qty{0.01}{s}$: $f = \qty{100}{Hz}$.
\end{solution}

\begin{solution}{exo:g9:alternating-current:3}
Cien veces: dos inversiones en cada uno de los cincuenta ciclos.
\end{solution}

\begin{solution}{exo:g9:alternating-current:4}
$T = 4 \times 5 = \qty{20}{ms} = \qty{0.02}{s}$; $f = 1/0.02 =
\qty{50}{Hz}$.
\end{solution}

\begin{solution}{exo:g9:alternating-current:5}
Los \qty{230}{V} son la \omterm{prop:g9:alternating-current:effective}{tensión eficaz} --- el valor continuo
constante que calentaría igual de bien, y lo que leen los
\omterm{def:g8:voltage-voltmeter:voltmeter}{voltímetros} de CA; los \qty{325}{V} son el pico, la cresta que
toca de verdad la oscilación, en cada sentido, cincuenta veces por
segundo.
\end{solution}

\begin{solution}{exo:g9:alternating-current:6}
Al calentamiento no le importa nada el sentido --- la \omterm{def:g8:ohms-law:resistance}{resistencia} se
calienta con cualquier desfile, así que los hervidores beben la
oscilación directamente. Cargar es química al revés, y la química exige
un sentido coherente: el ladrillo convierte primero la oscilación en
constancia.
\end{solution}

\begin{solution}{exo:g9:alternating-current:7}
La \omterm{def:g3:first-electric-circuit:battery}{pila}: una línea plana a \qty{9}{V}. La fuente alterna: una
onda suave oscilando entre $+9$ y \qty{-9}{V}, un ciclo cada
\qty{20}{ms}. La prueba del vistazo: plana frente a ondulada.
\end{solution}

\begin{solution}{exo:g9:alternating-current:8}
Las máquinas que giran fabrican CA de forma natural; y solo la CA
alimenta al transformador, sin el cual la electricidad no podría
recorrer el país a alta \omterm{def:g8:voltage-voltmeter:voltage}{tensión} y entrar en las casas con una
mansa.
\end{solution}

\begin{solution}{exo:g9:alternating-current:9}
Con pedaleo lento, los pocos ciclos por segundo de la dinamo hacen
pulsar la lámpara de manera visible --- cada cresta un destello, cada
cero un bajón; pedalear más rápido multiplica los pulsos hasta que el
ojo los funde. Las $24$ imágenes del cine fijan el ritmo aproximado de
\omterm{def:g6:water-states:changes}{fusión} del ojo: por debajo de unos \qty{25}{Hz} de pulsos de \omterm{def:g1:light-and-shadows:source}{luz}, el
parpadeo molesta; por encima, el resplandor se lee constante.
\end{solution}

\begin{solution}{exo:g9:alternating-current:10}
Crestas de unos \qty{325}{V} --- y cien por segundo, cincuenta
positivas y cincuenta negativas.
\end{solution}

\begin{solution}{exo:g9:alternating-current:11}
$T = 1/60 \approx \qty{0.017}{s}$; pico $\approx 120 \times 1.4 \approx
\qty{170}{V}$. Las estufas y las lámparas se adaptan con indiferencia;
los cargadores modernos llevan en su etiqueta ``100--240 V, 50/60 Hz''
y convierten tan contentos --- los que protestan son los aparatos
antiguos de una sola \omterm{def:g8:voltage-voltmeter:voltage}{tensión}, y cualquier reloj motorizado que
cuente los ciclos de la red: iría adelantado con sesenta.
\end{solution}

\begin{solution}{exo:g9:alternating-current:12}
Avivarse en las dos crestas da $2 \times 50 = 100$ pulsos por segundo.
Una cámara de $25$ fotogramas muestrea cada \qty{0.04}{s} --- cuatro
pulsos por fotograma, pero no alineados igual en cada punto de la
exposición que va barriendo: el pequeño desajuste entre el ritmo de la
cámara y el de la lámpara pinta bandas que se arrastran de un fotograma
al siguiente --- un batido entre dos relojes. A \qty{60}{Hz} ($120$
pulsos), el desajuste frente a los $25$ fotogramas cambia, y las barras
cambian de anchura y de \omterm{def:g7:motion-average-speed:formula}{velocidad} --- el dibujo delata la red
local.
\end{solution}

\begin{solution}{pb:g9:alternating-current:1}
\textbf{1.} CC a $2.25 \times 2 = \qty{4.5}{V}$: la \omterm{def:g3:first-electric-circuit:battery}{pila} plana de
bolsillo.
\textbf{2.} CA; $T = \qty{20}{ms}$, $f = \qty{50}{Hz}$; pico $3 \times
2 = \qty{6}{V}$.
\textbf{3.} CC a \qty{4.5}{V} conectada \emph{al revés} --- con las
puntas intercambiadas: inofensivo en un \omterm{met:g9:alternating-current:scope}{osciloscopio}, instructivo en el
registro.
\textbf{4.} CA; $T = \qty{10}{ms}$, $f = \qty{100}{Hz}$; pico
\qty{3}{V}. Sus pulsos doblados, cien por segundo, quedan más
cómodamente por encima del ritmo de \omterm{def:g6:water-states:changes}{fusión} del ojo que el parpadeo de
cincuenta ciclos de B.
\textbf{5.} El \omterm{def:g8:voltage-voltmeter:voltmeter}{voltímetro} declara el valor eficaz --- el
equivalente en continua para calentar ---, que para una sinusoide está
más o menos en el pico dividido entre $1.4$: $6 \div 1.4 \approx
\qty{4.3}{V}$. Los dos números describen una misma oscilación.
\textbf{6.} Solo la fuente A (con las puntas bien puestas --- y la C,
después de volver a enchufarla). Las fuentes de CA necesitan un
convertidor --- el ladrillo rectificador del cargador --- entre ellas y
cualquier batería.
\textbf{7.} Las cincuenta oscilaciones por segundo de la red están
cerca del propio ritmo eléctrico del corazón y pueden alterarlo:
\omterm{def:g8:intensity-ammeter:intensity}{amperio} a \omterm{def:g8:intensity-ammeter:intensity}{amperio}, la oscilación pone en peligro más
que el desfile constante.
\textbf{8.} La sinusoide --- la oscilación perfecta, dibujada de forma
natural por una bobina que gira a \omterm{def:g7:motion-average-speed:formula}{velocidad} constante en un campo
magnético: el alternador del capítulo siguiente.
\textbf{9.} Un ciclo cada $4$ cuadros (\qty{20}{ms}); crestas de $325
\div 100 = 3.25$ cuadros por encima y por debajo.
\textbf{10.} $f = 1/0.02 = \qty{50}{Hz}$; eficaz $\approx 325 \div 1.4
\approx \qty{230}{V}$ --- la etiqueta, recuperada de la pantalla.
\textbf{11.} Una línea plana a \qty{230}{V} --- la constancia de una
potencia calorífica igual.
\textbf{12.} Por ejemplo: ``Una línea plana es la constancia de una
\omterm{def:g3:first-electric-circuit:battery}{pila} --- un sentido, un valor. Una onda es la oscilación de la
red --- sentido y valor en ciclo incesante. Y el cambio justo entre las
dos es la \omterm{prop:g9:alternating-current:effective}{tensión eficaz}: mismo \omterm{def:g5:heat-and-insulation:heat}{calor}, mismo valor''.
\end{solution}
